% 2009 年浙江卷文科第 14 题：频率分布直方图。
% 来源 PDF：普通高考/2009/2009浙江文.pdf，SHA-256 8e609d19459cb4234a4d6dd9db5be5c8cc9af2169c00d4337b3988089cfb0709。
% 原页：第 1 页第 3 栏上部，书页码 567。
% 重绘依据：横轴 0—6，0 与 1 之间为省略断轴；纵轴为「频率/组距」。
%   各组矩形高：[1,2)=0.05，[2,3)=0.10，[3,4)=0.40，[4,5)=0.30，[5,6)=0.15。
%   虚线标出 0.05、0.10、0.15、0.40 四条水平参考线。
\begin{tikzpicture}[x=1.25cm, y=9.0cm, semithick,
  every node/.style={inner sep=1.5pt, font=\small}]

  % 坐标轴
  \draw[-{Stealth}] (-0.15,0) -- (6.55,0) node[right] {数据};
  \draw[-{Stealth}] (0,-0.025) -- (0,0.505);
  \node[align=center] at (0.58,0.468) {频率\\[-2pt]组距};

  % 水平虚线
  \foreach \h in {0.05,0.10,0.15,0.40} {
    \draw[dashed] (0,\h) -- (6.15,\h);
  }

  % 各组矩形
  \draw (1,0) rectangle (2,0.05);
  \draw (2,0) rectangle (3,0.10);
  \draw (3,0) rectangle (4,0.40);
  \draw (4,0) rectangle (5,0.30);
  \draw (5,0) rectangle (6,0.15);

  % 断轴符号（0 与 1 之间）
  \draw (0.22,0) -- (0.30,0.018) -- (0.38,0) -- (0.46,0.018) -- (0.54,0);

  % 刻度与标签
  \foreach \x in {1,2,3,4,5,6} {
    \draw (\x,0) -- (\x,-0.006);
    \node[below] at (\x,-0.008) {$\x$};
  }
  \node[below left=1pt] at (0,0) {$O$};
  \node[left] at (0,0.05) {$0.05$};
  \node[left] at (0,0.10) {$0.10$};
  \node[left] at (0,0.15) {$0.15$};
  \node[left] at (0,0.40) {$0.40$};
\end{tikzpicture}
